\documentclass[11pt,a4paper]{book}

\usepackage{bookmacro-lua-english}

\title{Introduction to Statistics}
\subtitle{Theory and Applications}
\author{Kazuma Tsuboi}
\date{\today}

\begin{document}

\maketitle
\tableofcontents

\chapter{Probability}

\section{Example Free Sentences}

“Order—arms!” said Major Tommy Toppleton, mounted on his sorrel pony, and facing the battalion formed of the students of the Toppleton Institute.

The butts of the light muskets clanged in unison on the gravel walk, where the column was drawn up. It was early on Monday morning, and the young soldiers were in excellent spirits—better than they were likely to be at a later hour in the day, for the programme included a fatiguing march. On the road leading from the Institute grounds stood three wagons, each drawn by two horses, upon which were loaded the baggage, provisions, and camp equipage of the battalion.

At the head of the column was the students’ band of ten pieces, besides a drum corps of eight, and two fifers. Signor Perelli, our music teacher, had bestowed incredible care and pains upon these amateur musicians, and those who were competent judges declared that the result was highly creditable to his skill and perseverance. In other words, the band played very well, though it did not undertake to compete with Gilmore, Strauss, Jullien, or even with the Ucayga Cornet Band. The drum corps was a perfect success, for most of its members had been in practice over a year.

Not a few of the people of Middleport had gathered in front of the Institute to witness the parade, for the occasion was no ordinary one, and the programme of the battalion had been thoroughly discussed for weeks. In years before the students had camped out during the June vacation. During the preceding season the Toppletonians had pitched their tents on the Horse Shoe, where the famous battles between the rival academies had taken place, as my friend Captain Wolf Penniman has related in his story.

Camping out had become rather stale with the students, and they longed for a new sensation. If they could have encamped on the Horse Shoe, and had another conflict with the Wimpletonians, perhaps it would have satisfied them. But there was no prospect of any sport of this kind, for Waddie Wimpleton, the haughty untamed, and tyrannical, was as mild and gentle as a lamb. It was said among the fellows that he had experienced religion, or something of that sort. It was certain he was not the boy he used to be. He had voluntarily resigned his positions as president of the Steamboat Company and major of the battalion. Though Ben Pinkerton was commander of the forces on the other side of the lake, it was supposed that Waddie’s influence as a peacemaker was sufficient to control the movements of the troops, and prevent them from engaging in another conflict with the fellows on our side of the lake.

Besides, there was another circumstance which seemed to interfere to keep the peace between the boys of the two Institutes. A military gentleman, residing at Ucayga, had become interested in the two battalions of juvenile soldiers, and had offered a prize of a magnificent standard to the one which should excel the other in company and battalion drill. With the vanity natural to boys, each party believed that all the skill and precision was upon its own side, and both had accepted the invitation to drill for the banner. This great event was to come off on the following Friday morning at Centreport. The place had been fixed by lot; and, when the arrangement had been completed, it suggested the present movement of our battalion.

The Toppletonians did not care to encamp at any point near enough to Centreport to enable them to keep their engagement, and when Captain Briscoe jocosely proposed that the little army should march round to the other side of the lake by the way of Hitaca, the idea was received with tremendous enthusiasm. The excursion would be a tour of camp duty, with an everchanging scene, and with no lack of novelty and excitement. We voted almost unanimously that the long tramp of seventy miles was just the thing we wanted.

It would afford us opportunity to display our new uniforms, our band and drum corps, our drill and marching, to people on the route who had hardly ever seen a company of soldiers. We should astonish the Hitacaites with our music and parade, excite the admiration of the ladies in general, and the young ladies in particular, and make us all first-class lions. The people would turn out to behold us, bestow delicate attentions upon us, and entertain us with generous hospitality. We had all the elements for a splendid parade, including our stylish uniforms, good music, and well-trained companies.

For two weeks hardly anything was talked about but this tour of camp duty. Those who received the plan coldly at first, soon became enthusiastic. Those who growled at the idea of marching twenty miles in a day were persuaded to believe that our progress would be a continued triumph, and that it would be accomplished in six or seven hours, so that there would be plenty of time to rest. The authorities of the Institute objected to the plan, but as Major Tommy Toppleton favored it, there was not much to be said against it. His father was compelled to indorse the march; and, of course, the instructors were obliged to withdraw all opposition.

Two of the teachers, besides the drill-master, were detailed to accompany the battalion in carriages, and prevent it from robbing hen-roosts, or capturing any of the towns on the route. But when Tommy heard of this little arrangement he was as indignant as though the professors had given him a thrashing, and interposed his veto. He would not have any schoolmasters dogging his steps, and spying into his actions. He hated spies. The battalion was composed of young gentlemen, and if they were a little fast at times, they knew how to behave themselves, and did not need any pedagogues to watch them. He could take care of his force himself. So the instructors were permitted to spend their vacation in pursuits more congenial to their tastes than following a multitude of crazy boys, under a crazy leader. Perhaps it would have been better for Tommy if he had permitted these guardians of the peace to attend the battalion.

We were all ready to start. The baggage wagons were to fall in behind the column, and the drivers were on their boxes. Everybody was in high spirits, and anticipated the greatest time known in the annals of the Toppleton Institute. Major Tommy, in particular, was in full feather; for he was the commander of the expedition, and a march of a week through places which suggested honors and ovations was an event which was calculated to stimulate his bump of self esteem.

On the preceding Saturday, the stockholders of the Lake Shore Railroad had held their annual meeting. On the year before, the present president had actually been defeated on the first ballot, and “your humble servant” elected in his place. Not caring to endure the constant browbeating and annoyance to which I should have been subjected had I taken the office, and because I really believed then that Tommy ought to have it, I had declined. Then, by a tremendous effort on the part of the president’s friends, and particularly on the part of Wolf Penniman, who pleaded for him as though he had been a friend in distress, Tommy received a bare majority of the votes cast.

I was sorry afterwards that I refused the position, for Tommy always overbearing and tyrannical, became to such an extent that the fellows had been on the verge of mutiny for a year. They had waited with impatience for the return of the annual elections, intending to pitch him down from his high positions. But Tommy and his father had provided for this emergency. The former had always nominally held—as, in fact, the rest of the students owned their stock—a large proportion of the shares. If a boy left the Institute during the year, his stock was made over to Tommy, or some of his toadies, for he had a small army of satellites of this species. The result was that the president was reëlected by a small majority. Feeling that he was secure in his high position, he had taken no pains to conciliate those who condemned his tyranny, and, if possible, he was more unpopular than ever.

The students had become pretty thoroughly disgusted with the management of the Lake Shore Railroad. Although the business of the company was still carried on in their name; though the stock was bought, sold and transferred on the books; though the boys discharged the duties of their several offices,—they had but little interest in the corporation. They had no real power. If the superintendent, the road-master, or the directors did anything, it was by order of Tommy or his father. But I ought to add, in justice to Major Toppleton, that there was not a student in the Institute, who had been there a year, that did not know all about the details of running a rail road, who was not familiar with stock operations, and who was not prepared to discharge his duty as an official in a railroad company. As a means of instruction it was still a good thing; but it had ceased to be a source of amusement, as it would have been if the president had not ruled so arbitrarily.

As the railroad president, Tommy was safe for another year. Though dissatisfied, our fellows had already begun to gather up the stock by helping out weaker ones in their lessons, or by the purchase of it with money, peanuts, and cream-cakes. The election of a commanding officer of the battalion usually came off on the same day as the railroad meeting; but it had been postponed to the following Tuesday, by what influence I do not know. Those who accomplished this purpose either forgot that on this day we should be on the march, or they expected to derive some advantage from the fact.

Tommy Toppleton declared that he was sure of being elected; and I think he was sincere in his belief. He had many devoted adherents, who bowed down to his power and influence; but there were just as many active opponents, and a great middle class who were not partisans either for or against him. Electioneering on both sides was carried on with spirit and energy, and the contest promised to be an exciting one.

\section{Random Variables}

A random variable assigns a numerical value to each outcome.

\begin{definition}[Random Variable]
A random variable is a measurable function from a sample space
to the real numbers.
\end{definition}

\begin{theorem}[Linearity of Expectation]
For random variables \(X\) and \(Y\),
\[
  \expectation[X+Y]
  =
  \expectation[X]+\expectation[Y].
\]
\end{theorem}

\begin{question}[Expectation]
Let \(X\) be a Bernoulli random variable. Find its expectation.
\end{question}

\begin{kaitou}
If $\mathrm{Pr}\paren{x=1} = p$, then
\[
  \expectation[X]=p.
\]
\end{kaitou}

\begin{defbox}{Variance}{def:variance}
The variance of \(X\) is
\[
  \operatorname{Var}(X)
  =
  \expectation\!\left[
    \left(X-\expectation[X]\right)^2
  \right].
\]
\end{defbox}

\begin{algorithm}
  \caption{Calculate the sample mean}
  \begin{algorithmic}
    \AlgInput Observations \(x_1,\dots,x_n\)
    \AlgOutput Sample mean \(\bar{x}\)
    \State \(\bar{x}\gets 0\)
    \For{\(i=1,\dots,n\)}
      \State \(\bar{x}\gets\bar{x}+x_i/n\)
    \EndFor
  \end{algorithmic}
\end{algorithm}


\appendix
\chapter{Additional Results}

Appendix material goes here.

\end{document}